\subsection{The filtration proposal of 119/153: a counterexample and an intrinsic repair}

The words ``regular equidimensional filtration'' on 119/153 do not
specify a frontier condition, a regularity test on closed unions, or
equisingularity. The original image, manuscript page~6, archive
foliation~152, confirms the correction to ``la moins fine'' and the
bracket concerning a preassigned filtration. We distinguish two precise
interpretations. A \emph{weak regular filtration} is a finite chain of
closed definable subsets whose nonempty successive differences are
manifolds without boundary of a single dimension. Its layers need not be
connected. Refinement means refinement of the associated layer
partitions; the counterexample below therefore applies also to the
more restrictive convention that a coarser chain is a subsequence.

\begin{saresult}{A.204}{Refuted}{There need not be a coarsest weak regular equidimensional filtration}
There is a compact semialgebraic subset $X\subset\R^4$ with two closed
filtrations, each having regular layers of dimensions $0,1,2$, which
have no common coarsening with regular layers. Thus the weak
interpretation of the coarsest-filtration claim at 119/153 is false,
both for embedded smooth regularity and for intrinsic PL regularity.
\end{saresult}
\begin{proof}
Let $N$ be the north pole of $S^2$, and identify $S^2\setminus\{N\}$
with $\R^2$ by stereographic coordinates $(x,y)$. Define
\[
 h(x,y)=
 \frac{\sqrt{(x-1)^2+y^2}+\sqrt{(x+1)^2+y^2}-2}
 {1+x^2+y^2},\qquad h(N)=0.
\]
The triangle inequality gives $h\ge0$, with equality in the finite
chart precisely on the segment $A=[-1,1]\times\{0\}$. If
$r=\sqrt{x^2+y^2}$, the numerator is at most $2r$ by the two
inequalities $|(x,y)-(\pm1,0)|\le r+1$. Thus
$0\le h\le2r/(1+r^2)$, proving continuity at $N$.
The function is semialgebraic and smooth outside the two endpoints
$a_-,a_+$ and $N$. Put
\[
 S_0=S^2\times\{0\},\qquad
 S_h=\{(z,h(z)):z\in S^2\},\qquad X=S_0\cup S_h.
\]
These are compact semialgebraic sets; both sheets are homeomorphic to
$S^2$, and $S_0\cap S_h=A\cup\{N\}$, with the customary
identification of their common points.

Let $C\subset S_0$ be the great circle given in the finite chart by
$y=0$, together with $N$. Consider the closed chains
\[
 \begin{array}{lll}
 F_0=\{a_-,a_+,N\},&F_1=A\cup\{N\},&F_2=X,\\
 G_0=\{N\},&G_1=C,&G_2=X.
 \end{array}
\]
Their respective layer partitions are
\[
 \begin{array}{lll}
 P_0=\{a_-,a_+,N\},&P_1=A\setminus\{a_-,a_+\},
 &P_2=X\setminus(A\cup\{N\}),\\
 Q_0=\{N\},&Q_1=C\setminus\{N\},&Q_2=X\setminus C.
 \end{array}
\]
The first layers are zero-dimensional smooth manifolds and the second
layers are smooth one-manifolds. The two sheets of $P_2$ are disjoint
open subsets of smooth graphs. Likewise $Q_2$ is the disjoint union
of the two open hemispheres of $S_0\setminus C$ and the smooth
graph $S_h\setminus(A\cup\{N\})$. Consequently $P_2$ and $Q_2$
are embedded smooth two-manifolds. A point of one of these pieces
has positive distance from the other sheet locally, so their union
has the asserted manifold structure even where the sheets approach
the removed sets. These facts also establish intrinsic PL regularity,
by semialgebraic triangulation of these smooth manifolds.

The join of the two partitions has a single atom. Indeed its
bipartite incidence graph contains the edges
\[
 P_0Q_0,\quad P_0Q_1,\quad P_1Q_1,\quad P_2Q_1,\quad P_2Q_2,
\]
which connect all six vertices. The edge $P_2Q_1$ is supplied by a
point of the base great circle outside $A\cup\{N\}$. Any partition
coarser than both must therefore identify all of $X$ into one layer.

The set $X$ is not a manifold. At an interior point of $A$, its germ
is two disks glued along a diameter, or four half-disks glued along
their common diameter. This description follows by using the sphere
projection separately on each sheet and identifying exactly their
shared diameter; it is a homeomorphism from the corresponding finite
gluing, irrespective of the tangency of the embedded sheets. A link
is the graph with two vertices joined by four edges. Its first
homology is $\Z^3$, so the local homology in degree two is
$\Z^3$. A two-manifold has local degree-two homology $\Z$.
Since $X$ has local dimension two, it cannot be a manifold of another
dimension either. Hence the sole common coarsening is not regular.
\end{proof}

The scope of this example matters. The partition $Q$ does not satisfy
the frontier condition after separating its connected components:
the closure of its upper-sheet two-stratum meets $Q_1$ precisely
along $A$, not along all of $Q_1$. It also fails a regularity test
on all closed saturated unions, because $\operatorname{Sing}(X)$
meets $Q_1$ in only a proper subset. We do not infer from A.204 that
the coarsest partition fails under either of those stronger readings.

\begin{saresult}{A.205}{Proved}{A canonical intrinsic PL partition and its dimension filtration}
Let $X$ be a definable set in a fixed o-minimal expansion
of a real closed field, and let $\Phi$ be a finite family of closed
definable subsets of $X$. There is a unique coarsest finite
partition $\mathcal P(X;\Phi)$ satisfying the following requirements:
its strata are locally closed PL manifolds, each of a single dimension;
each stratum closure is a union of strata; the sets in $\Phi$ are
unions of strata; and for every closed union $Y$ of strata, the
non-PL-manifold locus $\operatorname{Sing}_{\rm PL}(Y)$ is a union
of strata. Dimension may differ between strata. One may instead
require connected strata (definably connected over a general real
closed field); this class also has a unique coarsest
partition.

Both versions are preserved by definable homeomorphisms of the
marked data $(X;\Phi)$. Their dimension filtration
\[
 X_{\le k}=\bigcup_{\substack{S\in\mathcal P(X;\Phi)\\\dim S\le k}}S
 \quad(-1\le k\le\dim X)
\]
is closed and canonical, with pure-dimensional PL-manifold layers
after empty layers are omitted. If a closed filtration
$\varnothing=Y_{-1}\subset Y_0\subset\cdots\subset Y_s=X$ is
given, take $\Phi=\{Y_0,\ldots,Y_s\}$ and insert, in lexicographic
order, the closed sets
\[
 E_{j,k}=Y_{j-1}\cup(Y_j\cap X_{\le k}),
 \qquad 0\le j\le s,\quad -1\le k\le\dim X.
\]
Omitting repetitions gives a canonical regular equidimensional
filtration compatible with the preassigned one. The coarseness
assertion concerns the stated class of partitions; it is not the
false least-fine assertion among all weak filtrations.
\end{saresult}
\begin{proof}
We first justify the intrinsic operator, including for noncompact
sets. A finite family of definable sets admits a model by unions
of open faces of a finite compact polyhedron. To reduce to the
compact triangulation theorem, first make all sets bounded by the
coordinate homeomorphism $t\mapsto t/\sqrt{1+t^2}$. Take a finite
cell decomposition adapted to the sets. Each cell is locally
closed, and hence is the difference of its compact closure and
its compact frontier. Triangulate these finitely many closed
sets simultaneously inside their compact union. This gives the
required model from \cite[Theorem 2.2(1), p.~6]{Shiota2010}.
All closures used subsequently are relative to $X$.

Call a point intrinsically regular if it has a definable
neighborhood definably homeomorphic to an open Euclidean ball,
without boundary. The dimension is allowed to vary between local
components of a closed set; a regular isolated point of a
mixed-dimensional set is not declared singular. In any of the above polyhedral models this
is equivalent to PL regularity. Here is the local argument that
also proves independence of the model. Suppose a point $p$ is
intrinsically regular in a union of open faces of a finite
complex $K$. A chart supplies a compact definable ball $D$ whose
relative interior contains $p$ and whose center is $p$. Subdivide
$K$ so $p$ is a vertex. By
\cite[Theorem 2.2(2), p.~6]{Shiota2010}, a definable self-homeomorphism
$\tau$ of $|K|$ preserving each simplex, and therefore preserving
our union of open faces and fixing $p$, makes $\tau^{-1}(D)$
polyhedral. The relative Hauptvermutung
\cite[Theorem 2.1, p.~5]{Shiota2010}, applied with the centers
marked, gives a PL homeomorphism from a standard ball to this
polyhedron carrying its center to $p$. Since $\tau^{-1}(D)$
contains a relative neighborhood of $p$, the ambient germ is
PL Euclidean. The converse follows because a PL chart is definable.
Consequently this regularity is preserved by every definable
homeomorphism, even when the whole space is not locally compact.
Here compact over a non-Archimedean field means closed and bounded.

For a definable $Y$, the PL germ is constant on the relative
interior of each face in a finite model: it is the product of
that open face with its transverse conical germ. Its singular
locus is therefore a finite union of open faces and is definable.
The regular locus is open by definition, so the singular locus
is closed relative to $Y$. Its dimension is smaller than that
of $Y$, since interiors of the maximal included faces are regular
and dense. This verifies every regularity input needed below.

For completeness we give the finite existence and coarsening
argument with this intrinsic operator. Start with a compatible
finite triangulation of $(X;\Phi)$ and its open-face partition.
Freeze strata in decreasing dimension. Suppose strata of dimension
greater than $k$ are already frozen and the unfrozen remainder
is closed in $X$, partitioned with the frontier condition and
compatibly with the closures of all frozen strata. For every union
$A_J$ of closures of frozen strata and every current $k$-stratum
$S$, partition $S$ according to the finitely many sets
$S\cap\operatorname{Sing}_{\rm PL}(A_J)$. Remove their relative
boundaries and freeze the remaining open pieces, taking components
if the connected version is wanted. The removed boundaries have
dimension less than $k$. Every current $k$-stratum is open in the
unfrozen remainder, and hence so are the pieces just frozen. The
new remainder is therefore closed. Triangulate this remainder
compatibly with all frozen closures and the previous pieces. It
cannot accumulate onto a frozen stratum, because it is closed and
disjoint from the frozen strata. This maintains the frontier
condition and the invariant needed for the next dimension. Finiteness
of each stage and the dimension decrease make the process finite.

For a final $k$-stratum $S$ contained in a closed saturated $Y$,
the germ of $Y$ along $S$ is either the regular germ of $S$ alone
or the germ of a union $A_J$ of closures of higher-dimensional
strata of $Y$. Indeed the frontier condition makes each such
closure either contain all of $S$ or miss it, and finitely many
closures missing a given point can be avoided locally. In the
second case these higher strata had already been frozen when $S$
was selected, so the construction made
$\operatorname{Sing}_{\rm PL}(A_J)\cap S$ either empty or all of
$S$. Thus the singular locus of every closed saturated union is
saturated, not merely the singular locus of each stratum closure.

For uniqueness, write $\mathcal B_P$ for the finite Boolean
algebra of unions of strata of one such partition. This algebra is
closed under closure and under $Y\mapsto\operatorname{Sing}_{\rm PL}(Y)$
on closed members. For two admissible partitions $P,Q$, the
intersection $\mathcal B_P\cap\mathcal B_Q$ has those properties
and contains $\Phi$. Its atoms are locally closed. In fact, for an
atom $A$ its frontier $\overline A\setminus A$ and the closure of
that frontier belong to the algebra. If this closed frontier met
$A$, atomicity would imply that it contained $A$, contrary to
the strict frontier dimension inequality. Thus $A$ is open in
$\overline A$. The singular locus of $\overline A$ is saturated
and cannot contain the dense atom $A$, so $A$ is a PL manifold.
The atoms and their closures satisfy the frontier condition.

Purity is also preserved in this intersection. For every closed
saturated $Y$, the closed definable set
\[
 D_k(Y)=\{y\in Y:\dim_yY\ge k\}
\]
is saturated: its value on a stratum $S$ is determined by the
maximum dimension of a stratum whose closure contains $S$.
Thus both algebras, and their intersection, are closed under the
operators $D_k$. For an atom $A$, each $A\cap D_k(\overline A)$
is empty or all of $A$, so its manifold dimension is constant.
In the connected version the algebras are also closed under
individual connected components. A connected stratum lies in one
component of every saturated set containing it; components of such
sets are therefore saturated. This property passes to the algebra
intersection and forces its atoms to be connected.

Choose an admissible partition with the smallest possible number
of strata, which exists by the first construction and the
well-ordering of the positive integers. Its common coarsening with
any other admissible partition, constructed by the intersection
of algebras, cannot have fewer strata. It is therefore this same
partition, establishing coarseness and uniqueness. This is the
argument of A.303--304 with all regularity inputs verified for
the intrinsic PL operator.

A definable homeomorphism of marked data preserves closure, local
dimension, the PL regularity operator and connected components.
It therefore transports the unique coarsest partition to the
unique coarsest partition on the target, proving naturality.
The frontier of a stratum has strictly smaller dimension, so
$X_{\le k}$ contains the closure of every stratum it contains
and is closed. Distinct $k$-strata cannot accumulate onto one
another; their union $X_{\le k}\setminus X_{\le k-1}$ is therefore
a pure $k$-manifold. Finally, each $Y_j$ is saturated, so the
differences in the displayed refinement are exactly
\[
 (Y_j\setminus Y_{j-1})\cap
 (X_{\le k}\setminus X_{\le k-1}).
\]
They are unions of selected $k$-strata, hence pure $k$-manifolds
or empty. All sets and their order are canonically specified.
\end{proof}

The dimension filtration may merge strata with different incidence
or regularity roles. We do not assert that its merged layer
partition still satisfies the condition on all closed saturated
unions. Nor does that condition alone assert constant transverse
link type or the equisingularity required in the manuscript's
privileged-partition questions.

\begin{saresult}{A.206}{Refuted}{Embedded smooth regularity is not invariant under moderate homeomorphisms}
Replacing intrinsic PL regularity in A.205 by embedded smooth
regularity gives the semialgebraic coarsest partitions of A.303--304,
but does not give invariants under arbitrary definable homeomorphisms.
It does give naturality under isomorphisms preserving the embedded
smooth loci of every closed definable subset, for example ambient
smooth diffeomorphisms.
\end{saresult}
\begin{proof}
The map
\[
 [-1,1]\longrightarrow C=\{(t^2,t^3):-1\le t\le1\},
 \qquad t\longmapsto(t^2,t^3)
\]
is a semialgebraic homeomorphism; its inverse away from the origin
is $(x,y)\mapsto y/x$, extended continuously by zero. The point
$0$ is an embedded smooth point of the interval but its image is
not an embedded $C^1$ manifold point of the cusp. Indeed the tangent
cone there is one half-line, whereas a one-dimensional $C^1$
manifold without boundary has a full tangent line. Thus the smooth
singular-locus operator is not preserved. In the coarsest pure smooth
partition of the interval the open interval is one stratum, whereas
the cusp origin must be in a zero-dimensional stratum, because
the singular locus of the whole compact cusp is saturated. No
homeomorphism between these marked spaces can identify these two
partitions. The positive assertion follows by transporting all
the defining conditions, just as in the naturality proof of A.205.
\end{proof}

\paragraph{Precisely determined d\'evissage invariants.}
For the intrinsic construction the finite frontier poset, defined by
$S\le T$ if $S\subset\overline T$, with its dimension labels and
its membership labels in the original filtration, is an invariant.
The definable homeomorphism types of the strata and the filtered
definable homeomorphism type of $X$ are invariants as well. When
$X$ is compact they determine invariant PL types: apply relative
Hauptvermutung to the finite family of stratum closures and their
frontiers. No global noncompact Hauptvermutung is needed or asserted.
These assertions
mean preservation under marked definable homeomorphisms; they do
not assert a finite numerical classification of tame spaces.

For compact $X$ and a coefficient ring $\Lambda$ the finite dimension filtration
also determines the full cohomological spectral sequence
\[
 \begin{gathered}
 E_1^{p,q}=H^{p+q}(X_{\le p},X_{\le p-1};\Lambda)
 =H_c^{p+q}(X_{\le p}\setminus X_{\le p-1};\Lambda),\\
 E_1^{p,q}\ \Longrightarrow\ H^{p+q}(X;\Lambda),
 \qquad d_1:E_1^{p,q}\longrightarrow E_1^{p+1,q}.
 \end{gathered}
\]
Choose a compatible finite triangulation of the closed filtration.
The relative simplicial cochain complexes and their short exact
sequences give its exact couple, hence the displayed sequence;
the bounded filtration proves convergence. Naturality follows
from the functoriality of relative cohomology. Over a general real
closed field these groups mean those of a finite abstract
polyhedral realization, equivalently definable cohomology; relative
triangulation and Hauptvermutung make the construction independent
of the model. The same construction applies to the inserted
filtration $E_{j,k}$. In particular, when the coefficient dimensions
are finite,
\[
 \chi(X)=\sum_{S\in\mathcal P(X;\Phi)}\chi_c(S).
\]
This identity follows by additivity for each closed filtration
step, or by counting open simplices. The connecting maps and
higher differentials are part of the invariant: the list of Betti
numbers of the strata by itself does not supply them.

\subsection{A fixed-tube obstruction to numerical uniqueness}

\begin{saresult}{A.207}{Refuted}{The dimension condition leaves a full concordance space of rigidifications}
Let $Y$ be a closed compact definable manifold, let
$X=Y\times[0,2]$, and fix the strict tube
$V=Y\times[0,1]$, with core $Y\times\{0\}$ and identified outer
face $B=Y\times\{1\}$. The simplicial space of $Y$-admissible
rigidifications of this fixed tube satisfying
\[
 \dim\varphi^{-1}(y)=\operatorname{codim}_y(Y,X)-1=0
 \quad\text{for every }y\in Y
\]
is naturally isomorphic, after interval reversal, to
$P_{\rm def}(Y)$. Consequently the proposed dimension condition
fails to give essential uniqueness even when the tube and its
outer-face identification are fixed. For $Y=(S^1)^6$ this space
has the nonzero component group of A.103.
\end{saresult}
\begin{proof}
Throughout this argument over an arbitrary real closed field,
compactness means definable compactness and connectedness means
definable connectedness. The separations used below are definable;
closed-image and half-ball arguments are valid in that sense.
An admissible rigidification supplies a continuous map
\[
 q:B\times[0,1]\longrightarrow V,
 \qquad q(b,1)=b,\quad q(b,0)=(\varphi(b),0),
\]
which is a homeomorphism from $B\times(0,1]$ onto $V\setminus Y$
and, together with the identity on $Y$, induces the specified
mapping-cylinder homeomorphism. We first show that every fibre
of $\varphi$ is connected. This argument also explains exactly
why finite fibres force $q$ itself to be a homeomorphism.

Surjectivity of $\varphi$ is necessary. Its image is closed by
compactness. If $y$ were outside that image, a relative open
neighborhood of $y$ in the mapping cylinder would consist only
of points of $Y$. This contradicts the boundary half-space chart
of $V$ at $y$.

Suppose a fibre has a separation into two nonempty disjoint
compact sets $F_1,F_2$. Choose disjoint open neighborhoods
$O_1,O_2$ of these sets in $B$. The compact set
$B\setminus(O_1\cup O_2)$ has image avoiding $y$, so there is an
open neighborhood $U$ of $y$ in $Y$ such that
$\varphi^{-1}(U)\subset O_1\cup O_2$. For $0<\epsilon<1$, the
set
\[
 N=U\ \cup\ q\bigl(\varphi^{-1}(U)\times[0,\epsilon)\bigr)
\]
is open in the mapping cylinder and hence in $V$. Its complement
of the core is partitioned into the two disjoint relatively open
and closed pieces coming from
$\varphi^{-1}(U)\cap O_i$, times $(0,\epsilon)$.
Choose a boundary half-ball neighborhood $W$ of $y$ in $V$
contained in $N$. Then $W\setminus Y$ is connected. But for
$b_i\in F_i$ the rays $q(b_i,r)$ tend to $y$ as $r$ tends to
zero, so $W\setminus Y$ meets both pieces. This is impossible.
Thus each fibre is connected.

Under the dimension condition every fibre is finite, by
zero-dimensional definable finiteness, and is therefore a single
point. The map $\varphi:B\to Y$ is a continuous definable
bijection of compact sets, hence a homeomorphism. The displayed
map $q:B\times[0,1]\to V$ is now also a continuous definable
bijection, hence a homeomorphism. Conversely every homeomorphism
of the cylinder fixed pointwise on its top is admissible: its
bottom restriction is a homeomorphism $\varphi$, it identifies
$M_\varphi$ with $V$, and its fibres satisfy the dimension
condition. Boundary invariance ensures that it preserves the
bottom setwise. This construction and its inverse preserve all
parameter families, and interval reversal identifies the result
with $P_{\rm def}(Y)$.

This is a fixed-tube statement and does not use variation of the
neighborhood itself. For the torus, A.103 gives nontrivial
components and hence disproves contractibility and essential
uniqueness in this numerical class.
\end{proof}

\paragraph{The unrestricted admissible-rigidification gate.}
Dropping the dimension condition enlarges the space in A.207.
The proof still forces connected fibres in the product-tube
example, but such fibres can be nontrivial continua. In
particular it is invalid to identify every admissible
rigidification with a homeomorphism of the uncollapsed cylinder.
For example a PL monotone map of a polygonal circle that collapses
one closed edge to a vertex and maps its remaining edges once
around the circle has an annulus as its mapping cylinder: at
the collapsed boundary vertex its link is an interval, and all
other points have the usual annulus charts. Its attachment map
is not injective. Thus nontriviality of the homeomorphic-attachment subspace by
itself does not prove nontriviality of the entire admissible space.
The additional argument in A.209 below supplies this missing step.

Here is an exact finite certificate for the annulus assertion.
Let the source polygon have cyclic vertices $b_0,b_1,b_2,b_3$,
let the target have cyclic vertices $c_0,c_1,c_2$, and define
$\varphi(b_0)=\varphi(b_1)=c_0$, $\varphi(b_2)=c_1$,
$\varphi(b_3)=c_2$, linearly on each edge. Its mapping cylinder
has a triangulated model with the following seven oriented
triangles:
\[
\begin{gathered}
(b_0,b_1,c_0),\quad (b_1,b_2,c_0),\quad
(b_2,c_1,c_0),\quad (b_2,b_3,c_1),\\
(b_3,c_2,c_1),\quad (b_3,b_0,c_2),\quad
(b_0,c_0,c_2).
\end{gathered}
\]
Each uncollapsed edge contributes its usual triangulated square.
The collapsed-edge square contributes the single triangle
$(b_0,b_1,c_0)$: its quotient parametrization is
$(u,r)\mapsto (1-r)c_0+r((1-u)b_0+ub_1)$.
This is semialgebraic, is injective for $r>0$, and identifies
exactly the bottom edge at $r=0$. These parametrizations agree
on the vertical edges, proving that the displayed complex is
the actual mapping cylinder. No assertion that the quotient
parametrization is PL is needed.

The complex has $7$ vertices, $14$ edges, and $7$ triangles.
Every vertex link is an interval, each edge belongs to one or
two triangles, the listed orientations cancel on every interior
edge, and its boundary is exactly the disjoint source four-cycle
and target three-cycle. Thus it is a connected orientable compact
surface with two boundary components and Euler characteristic
zero, hence an annulus. The attached script
\texttt{check\_circle\_mapping\_cylinder.py} verifies all these
finite incidences and link conditions, as well as
$\partial_1\partial_2=0$ and the exact rational boundary ranks
$(6,7)$; all $17$ assertions pass. Consequently the target
boundary is collared while its attachment fibre over $c_0$ is a
nontrivial closed interval.

For a specified tube $(V,Y,B)$, the unrestricted question concerns
all maps $q$ as above, with jointly continuous parameter families,
including singular attachment fibres. Homotopy uniqueness of
$\varphi$ from A.201 does not imply contractibility of this space.
Nor does the point-core Alexander contraction apply: shrinking toward
a positive-dimensional core need not shrink tangential motion.
A.208 identifies the singular fibres, and A.209 then proves that the
concordance obstruction survives their addition.

\subsection{Contractibility of attachment fibres at a boundary core}

\begin{saresult}{A.208}{Proved}{Every admissible attachment fibre at a manifold boundary is contractible}
Let $V$ be a compact definable manifold with boundary, let $Y$ be a
union of boundary components, and let $B$ be a compact definable set.
Suppose a continuous definable map
\[
 q:B\times[0,1]\longrightarrow V
\]
has $q(b,0)=\varphi(b)\in Y$ and restricts to a homeomorphism
$B\times(0,1]\to V\setminus Y$. Then $\varphi$ is onto, and every
fibre $\varphi^{-1}(y)$ is definably contractible.

Over $\mathbb R$, the same neighborhood argument without definability
shows that all fibres are cell-like when $B$ is a compact ANR. Thus
all unrestricted attachment maps for the product tube in A.207 are
cell-like, and their definable fibres are actual contractible compact
polyhedra up to definable homeomorphism.
\end{saresult}
\begin{proof}
All compactness and connectedness arguments in the arbitrary-field
case are definable. The image of $q$ is closed, because its domain is
definably compact, and contains the dense subset $V\setminus Y$.
It is therefore all of $V$, and its bottom map $\varphi$ is onto.
The induced map $M_\varphi\to V$ is a continuous definable bijection
of definably compact spaces and hence a homeomorphism.

Fix $y\in Y$, put $F=\varphi^{-1}(y)$, and let $O$ be any definable
open neighborhood of $F$ in $B$. The compact set $B\setminus O$ has
closed image avoiding $y$. There is consequently an open neighborhood
$U\subset Y$ of $y$ such that $\varphi^{-1}(U)\subset O$.
For $0<\epsilon<1$ the subset
\[
 N=U\ \cup\ q\bigl(\varphi^{-1}(U)\times[0,\epsilon)\bigr)
\]
is open in the mapping cylinder, hence in $V$. Take a definable
boundary half-ball $W$ about $y$ contained in $N$. Its complement
$W\setminus Y$ is definably homeomorphic to a convex open half-ball,
so is definably contractible.

Continuity of $q$ and definable compactness of $F$ give some
$0<\delta<\epsilon$ with
$q(F\times\{\delta\})\subset W\setminus Y$. For completeness,
if no such uniform $\delta$ existed, definable choice would give
$b(r)\in F$ with $q(b(r),r)\notin W$ for arbitrarily small positive
$r$. Curve selection and compactness give a limiting $b_0\in F$,
contradicting $q(b_0,0)=y\in W$ and continuity. Contract
$q(F\times\{\delta\})$ to a point in $W\setminus Y$, and compose
this contraction with
\[
 W\setminus Y\ \xrightarrow{\ q^{-1}\ }\
 \varphi^{-1}(U)\times(0,\epsilon)
 \xrightarrow{\operatorname{pr}_B}\varphi^{-1}(U)\subset O.
\]
At time zero this is the inclusion $F\hookrightarrow O$.
Thus $F$ contracts in every such neighborhood $O$.

To obtain a contraction in $F$ itself, triangulate the compact pair
$(B,F)$ and take one barycentric subdivision. Write it as $(|K|,|L|)$,
with $L$ a full subcomplex. If $a_v(x)$ are the barycentric coordinates,
put
\[
 O_L=\left\{x:\sum_{v\in L}a_v(x)>0\right\},\qquad
 r(x)=\frac{\sum_{v\in L}a_v(x)v}{\sum_{v\in L}a_v(x)}.
\]
Fullness makes $r(x)\in|L|$, and this is a continuous semialgebraic
retraction $O_L\to|L|$. The segment from $x$ to $r(x)$ stays in its
simplex. Transport this neighborhood and retraction to $(B,F)$,
apply the preceding contraction into that neighborhood, and compose
with the retraction. The result is a definable contraction of $F$.

In the real topological setting ordinary compactness supplies the
same $U$, $W$, and uniform $\delta$ for every open $O$. It follows
that the inclusion of $F$ into every neighborhood in the ANR $B$
is nullhomotopic, which is the cell-like criterion. Definability is
used only for the stronger last step that contracts inside $F$.
\end{proof}

\subsection{Unrestricted admissible rigidifications retain distinct components}

\begin{saresult}{A.209}{Refuted}{Essential uniqueness fails even when all attachment fibres are allowed}
Let $Y=(S^1)^6$, let $V=Y\times[0,1]$, and fix its core
$Y\times\{0\}$ and outer face $B=Y\times\{1\}$.
Let $\mathcal R$ be the space of all continuous maps
$q:Y\times[0,1]\to V$ satisfying
\[
 q(y,1)=(y,1),\qquad q(Y\times\{0\})\subset Y\times\{0\},
 \qquad q|_{Y\times(0,1]}:Y\times(0,1]\xrightarrow{\sim}
                      V\setminus(Y\times\{0\}).
\]
Give it the compact-open topology. Its homeomorphic-attachment
subspace is $G=\Homeo(V\ \mathrm{rel}\ B)$, the top-fixed
concordance space. There is a locally constant map
\[
 \lambda:\mathcal R\longrightarrow\pi_0G
 \quad\text{with}\quad\lambda(g)=[g]\quad(g\in G),
\]
where $\pi_0G$ is discrete. In particular
$\pi_0G\to\pi_0\mathcal R$ is injective, so $\mathcal R$ is not
contractible and not even path connected.

The same counterexample applies to jointly continuous definable
families in any o-minimal expansion of $\mathbb R$. It also applies
to semialgebraic families over every real closed field, with a fixed
finite PL representative of the two distinct rigidifications.
Thus essential uniqueness in the full fixed-tube sense fails; no
fibre-dimension restriction is needed for this failure. The tube is
a strict tube in $X=Y\times[0,2]$.
\end{saresult}
\begin{proof}
We first give the uniform separation argument. For any compact
metric space $Z$ and subgroup $G\subset\Homeo(Z)$, set
\[
 d_\infty(g,h)=\sup_{z\in Z}d(gz,hz).
\]
This metric induces the compact-open group topology and is right
invariant:
\[
 d_\infty(gh^{-1},\id)=d_\infty(g,h).
\]
Suppose $G$ is locally path connected at the identity. There is
$\epsilon>0$ such that every $g\in G$ with
$d_\infty(g,\id)<\epsilon$ lies in the identity path component.
Consequently two elements less than $\epsilon$ apart lie in the
same path component. Distinct components are uniformly separated.

For any $f$ in the closure $\overline G$ in $C(Z,Z)$, choose
$g\in G$ with $d_\infty(f,g)<\epsilon/4$, and define
$\lambda(f)=[g]$. Any two such choices are less than
$\epsilon/2$ apart, proving independence. If
$d_\infty(f,f')<\epsilon/4$, approximations of $f,f'$ within
$\epsilon/4$ are less than $3\epsilon/4$ apart; hence their
labels agree. Thus $\lambda:\overline G\to\pi_0G$ is locally
constant and restricts to the usual component map on $G$.
This argument uses no continuous selection of approximations.

For $Z=V$ the required local path connectedness of the top-fixed
subgroup follows from the local contractibility of compact-manifold
homeomorphism groups
\cite[Corollary~1.1, p.~63]{EdwardsKirby1971}.
The subgroup preserving the two boundary components is open. If
$h$ belongs to it, write $a=h|_B:Y\to Y$ and define
\[
 R(h)=(a^{-1}\times\id_{[0,1]})h.
\]
This is a continuous retraction to $G$ and fixes the identity.
Restricting a sufficiently small local contraction and then applying
$R$ proves the local contractibility of $G$. This is also the
retraction used for the shell group in A.117.

It remains to prove $\mathcal R\subset\overline G$. Let
$q\in\mathcal R$. The real topological part of A.208 says that
its bottom attachment has cell-like fibres. Every other fibre of
$q$ is a point. Compactness makes $q$ proper and onto. Its
restriction to the whole boundary is also cell-like: it is the
bottom attachment on one component and the identity on the other.
Thus the approximation theorem
\cite[Theorem~A, p.~271]{Siebenmann1972Approximation} applies to
$q:V\to V$. The dimension is seven, so its case (a),
$m\ne3,4,5$, applies without an exceptional boundary hypothesis.
For every $\delta>0$ it gives a homeomorphism $h:V\to V$
with $d_\infty(h,q)<\delta$.

Use a product metric on $V$ and choose $\delta$ smaller than half
the distance between the two boundary components. Since $q$ fixes
the top, $h$ preserves that boundary component. Put
$a=h|_B$ and $g=(a^{-1}\times\id)h$. Then $g\in G$.
Moreover $d_\infty(a,\id)<\delta$ on $Y$, and
$d_\infty(a^{-1},\id)=d_\infty(a,\id)$. Hence
\[
 d_\infty(g,q)\leq d_\infty(g,h)+d_\infty(h,q)<2\delta.
\]
This proves the required density. The component label constructed
above now restricts to $\mathcal R$. A continuous path in
$\mathcal R$ has constant label, proving the asserted injection.

By A.103 and A.113 there is a PL top-fixed concordance
$p:V\to V$ with $[p]\ne[\id]$ in $\pi_0G$.
Both $p$ and $\id$ are admissible rigidifications of this same tube:
for $p$ the bottom attachment is $p|_{Y\times\{0\}}$, and the
induced mapping-cylinder identification is the identity on the
actual target core. The label prevents any path between them,
even if the path allows every singular attachment admitted by
$\mathcal R$. The exterior product $Y\times[1,2]$ supplies the
strictness collar. A jointly continuous real definable family is a
continuous path in $\mathcal R$, so it cannot connect these two
objects either. Any finite simplicial path gives such a path after
concatenation.

For the semialgebraic statement, A.115 permits $Y$ and $p$ to be
chosen over the real algebraic numbers. A proposed semialgebraic
path over a real closed extension is specified by finitely many
graph formulas and coefficient tuples. Joint continuity, the
outer-face identity, the bottom-face condition, and fibrewise
bijectivity off the bottom are first-order conditions, as are the
two prescribed endpoint formulas. These conditions transfer to the
real algebraic numbers and then to $\mathbb R$. Each transferred
map off the bottom is indeed a homeomorphism: the total compact
map is a quotient, and its restriction over the open complement of
the core is a continuous bijection with continuous inverse.
The transferred family would therefore be a real path in
$\mathcal R$ from $\id$ to $p$, which has just been excluded.
This step does not assume a real model for a general o-minimal
expansion of a non-Archimedean field.
\end{proof}

The proof deliberately uses only the individual approximation theorem
and an explicit uniform component-separation argument. The paragraph
after the complement on p.~272 of
\cite{Siebenmann1972Approximation} defers the nearly canonical and
parameter-dependent version to another manuscript. Neither that
unpublished version nor a deformation preserving full rulings is
needed here. A.209 refutes essential uniqueness; it does not assert
that the inclusion $G\hookrightarrow\mathcal R$ is a weak equivalence
or compute every component of the larger space.
